% !TeX root = ../navier-stokes-manuscript.tex
% ============================================================
% 1.7 Singularity and the Derivative Tower
% ============================================================
\subsection{Singularity and the Derivative Tower}\label{sub:SingularityAndTheDerivativeTower}

A singularity does not need to announce itself by sending the velocity to infinity.
The velocity may remain finite while its gradient becomes unbounded, while two approaches produce different gradients, or while a higher derivative stops existing.
Smoothness asks for every finite derivative order, so the relevant object is the entire response stacked above \(u\).

\divider{From Velocity to Acceleration and Jerk}

The first response is already present in the equation.
For the material derivative
\[
  D_t:=\partial_t+u\cdot\nabla,
\]
the acceleration of a moving fluid particle is
\[
  a:=D_tu=-\nabla p+\nu\Delta u.
\]
The left side says how the particle velocity changes.
The right side says that the change is determined by the pressure and viscous response of the field containing that particle.

Following the same particle once more gives the material jerk
\[
  j:=D_ta=D_t^2u.
\]
This derivative is different from the Eulerian second time derivative \(\partial_t^2u\), because the particle changes position as time passes.
If
\[
  V_n:=\partial_t^n u,
\]
then direct expansion gives
\[
  D_t^2u
  =
  V_2
  +2(V_0\cdot\nabla)V_1
  +(V_1\cdot\nabla)V_0
  +(V_0\cdot\nabla)\bigl((V_0\cdot\nabla)V_0\bigr).
\]
Acceleration and jerk bring spatial derivatives into the temporal response before any separate hierarchy has been invented.

The Eulerian hierarchy makes that nesting easier to see.
Set
\[
  P_n:=\partial_t^n p.
\]
The original equation gives the first temporal rung
\[
  V_1
  =
  \nu\Delta V_0
  -(V_0\cdot\nabla)V_0
  -\nabla P_0.
\]
Differentiating once gives
\[
  V_2
  =
  \nu\Delta V_1
  -(V_1\cdot\nabla)V_0
  -(V_0\cdot\nabla)V_1
  -\nabla P_1.
\]
Differentiating again gives
\[
  V_3
  =
  \nu\Delta V_2
  -(V_2\cdot\nabla)V_0
  -2(V_1\cdot\nabla)V_1
  -(V_0\cdot\nabla)V_2
  -\nabla P_2.
\]
The next rung is
\[
  \begin{aligned}
  V_4
  ={}&
  \nu\Delta V_3
  -(V_3\cdot\nabla)V_0
  -3(V_2\cdot\nabla)V_1\\
  &-3(V_1\cdot\nabla)V_2
  -(V_0\cdot\nabla)V_3
  -\nabla P_3.
  \end{aligned}
\]
The coefficients \(1,1\), then \(1,2,1\), then \(1,3,3,1\), are the rows of Pascal's triangle.
They record every way the time derivatives can split across the two velocity factors in the transport term.

At arbitrary order, the same calculation becomes
\[
  V_{n+1}
  +
  \sum_{k=0}^{n}
  \binom{n}{k}
  (V_k\cdot\nabla)V_{n-k}
  +
  \nabla P_n
  =
  \nu\Delta V_n,
  \qquad
  \nabla\cdot V_n=0.
\]
This recurrence is the temporal spine of the tower.
Every new rung is tied to the lower rungs through transport, to two additional spatial derivatives through viscosity, and to the pressure response at the same order.

\divider{Pressure and the Spatial Rungs}

Pressure is nested into the hierarchy as well.
Taking the divergence of the original momentum equation and using \(\nabla\cdot u=0\) gives
\[
  -\Delta p
  =
  \partial_i\partial_j(u_i u_j).
\]
Here and below, repeated spatial indices are summed.
After \(n\) time derivatives, the pressure at the \(n\)-th rung satisfies
\[
  -\Delta P_n
  =
  \partial_i\partial_j
  \sum_{k=0}^{n}
  \binom{n}{k}
  V_{k,i}V_{n-k,j}.
\]
Thus \(P_n\) is determined by the velocity rungs through an elliptic problem over the domain.
The pressure at one location generally depends on the configuration of the velocity field far beyond an arbitrarily small neighbourhood of that location.

Spatial differentiation opens the other direction of the tower.
For one coordinate direction \(x_\ell\), differentiating the momentum equation gives
\[
  \begin{aligned}
  \partial_t(\partial_\ell u)
  &+
  (u\cdot\nabla)(\partial_\ell u)
  +
  ((\partial_\ell u)\cdot\nabla)u
  +
  \nabla(\partial_\ell p)\\
  &=
  \nu\Delta(\partial_\ell u).
  \end{aligned}
\]
At the next spatial rung, two differentiated velocity factors appear:
\[
  \begin{aligned}
  \partial_t(\partial_{\ell m}u)
  &+
  (u\cdot\nabla)(\partial_{\ell m}u)
  +
  ((\partial_\ell u)\cdot\nabla)(\partial_m u)\\
  &+
  ((\partial_m u)\cdot\nabla)(\partial_\ell u)
  +
  ((\partial_{\ell m}u)\cdot\nabla)u
  +
  \nabla(\partial_{\ell m}p)\\
  &=
  \nu\Delta(\partial_{\ell m}u).
  \end{aligned}
\]
Repeated time and spatial differentiation continues the same product-rule pattern.
Using a multi-index \(\alpha=(\alpha_1,\alpha_2,\alpha_3)\) for spatial derivatives, collect the resulting velocity values in
\[
  \mathcal J^\infty u(t,x)
  :=
  \left\{
    \partial_t^n\partial_x^\alpha u(t,x):
    n\geq 0,\ 
    \alpha\in\mathbb N_0^3
  \right\}.
\]
This set is the full derivative tower at \((t,x)\).
Its base is the velocity, its temporal rungs contain the Eulerian responses \(V_n\), and its spatial directions contain the gradients and curvatures on which transport and viscosity act.
The pressure tower is coupled to it through the differentiated Poisson equations.

No rule requires every rung to be nonzero.
The zero solution has a completely vanishing tower, and a steady solution has \(V_n=0\) for every \(n\geq1\), while both remain perfectly smooth.
The word \emph{live} must describe lawful participation rather than positive size.
A live tower has finite, compatible values that remain derivatives of one field and satisfy the differentiated equations wherever those equations are claimed.

This is the useful singularity picture.
A proposed endpoint can fail first at velocity, acceleration, jerk, a spatial derivative, or any mixed rung above them.
Whatever the first visible rung may be, the failure belongs to one tower whose levels are already nested into the Navier--Stokes law.
